% Image credits for the photographs and AI illustrations of Book 4
% (University Physics, Year 2). Machine-readable license records live in
% images/book4/CREDITS.md; keep the two files in sync.
\chapter*{Créditos de las imágenes}

Los dibujos de este libro son originales. Las fotografías se usan bajo
sus respectivas licencias libres, con agradecimiento a sus autores:

\begin{itemize}
  \item \emph{Tuberías forzadas de una central hidroeléctrica}
        (capítulo 5): Chinchu.c, Wikimedia Commons, CC0.
  \item \emph{Calle de torbellinos de von Kármán en las nubes}
        (capítulo 4): Robert Cahalan, NASA/GSFC, dominio público.
  \item \emph{Cable coaxial} (capítulo 8): Mx.~Granger, Wikimedia
        Commons, CC0.
  \item \emph{James Clerk Maxwell} (capítulo 11): Smithsonian
        Institution, dominio público.
  \item \emph{Aurora austral vista desde la Estación Espacial
        Internacional} (capítulo 14): Joe Acaba, NASA, dominio
        público.
  \item \emph{Franjas de interferencia entre dos superficies planas}
        (capítulo 19): revista \emph{Machinery}, hacia 1915, dominio
        público.
  \item \emph{LIGO Hanford, vista aérea} (capítulo 20): LIGO
        Laboratory, dominio público.
  \item \emph{Láser de helio--neón a 594~nm} (capítulo 23):
        Telementor, Wikimedia Commons, CC~BY~4.0.
  \item \emph{El Sol en el ultravioleta extremo} (capítulo 26):
        NASA/SDO (AIA), dominio público.
  \item \emph{La Tierra fotografiada desde el Apolo 17}
        (capítulo 26): tripulación del Apolo 17, NASA, dominio
        público.
  \item \emph{Ludwig Boltzmann} (capítulo 29): fotógrafo desconocido,
        dominio público.
  \item \emph{Figura de difracción de electrones} (capítulo 30):
        National Institute of Standards and Technology, dominio
        público.
  \item \emph{Superficie de oro (100) vista con el microscopio de
        efecto túnel} (capítulo 31): Erwin Rossen, Wikimedia Commons,
        dominio público.
  \item \emph{Charles Townes y el primer máser} (capítulo 31): Dan
        Rubin, dominio público.
\end{itemize}

Los enlaces a las fuentes completas y las direcciones de las licencias
de cada fotografía figuran en \texttt{images/book4/CREDITS.md}, en el
repositorio del libro.

\bigskip
Junto a las fotografías, treinta figuras de los capítulos 1--10 y
12--31 son ilustraciones generadas por inteligencia artificial,
producidas con la generación de imágenes de OpenAI a partir de
indicaciones escritas por los editores del libro, y revisadas para
comprobar su exactitud física antes de incluirlas. La indicación
completa de cada imagen generada está recogida en
\texttt{images/book4/ai/PROMPTS.md}, en el repositorio.
\clearpage
